\section{Grammar, notation and history spaces}\label{sec:grammar}
\subsection{The grammar}
\begin{definition}[SELECT-II]\label{def:select}
For an integer \emph{label} $N\ge1$ put $R(N)=(N+1)^2/N$ and $T_N(m)=m\,R(N)$. The label $N$ is \emph{admissible at} the integer $m\ge1$ if $T_N(m)\in\mathbb Z$.
\end{definition}
\begin{lemma}[admissibility is divisibility]\label{lem:adm}
For a positive integer $m$, $T_N(m)\in\mathbb Z$ if and only if $N\mid m$.
\end{lemma}
\begin{proof}
$\gcd(N,(N+1)^2)=1$, so $N\mid m(N+1)^2$ iff $N\mid m$.
\end{proof}
\begin{lemma}[multipliers]\label{lem:R}
$R(N)=N+2+1/N$; $R(1)=4$, $R$ is strictly increasing on integers $N\ge1$, and hence $R(N)\ge4$ with equality only for $N=1$. In particular $N\mapsto R(N)$ is injective.
\end{lemma}
\begin{proof}
$R(N+1)-R(N)=1-\frac1{N(N+1)}>0$.
\end{proof}
For a prime $p$ let $v_p$ be the $p$-adic valuation and put
\[
 \Delta_p(N)=v_p\bigl(R(N)\bigr)=2v_p(N+1)-v_p(N),
\]
so at most one of $v_p(N),v_p(N+1)$ is nonzero and $\Delta_p(N)\ge-v_p(N)$ with equality unless $p\mid N+1$. A prime is \emph{active} for a set of labels if it divides one of them. All valuations of $R(N)$ vanish outside the primes dividing $N(N+1)$.

\begin{definition}[history]\label{def:history}
A \emph{history} is a pair $h=(m_0;w)$ with $m_0\in\mathbb Z_{\ge1}$ the \emph{source} and $w=(N_1,\dots,N_l)$ a word of labels such that $m_j=T_{N_j}(m_{j-1})$ is an integer for $j=1,\dots,l$, i.e.\ $N_j\mid m_{j-1}$. The integers $m_0,\dots,m_l$ are the \emph{states} of $h$, $m_l$ is the \emph{endpoint}, $l$ the length, and the multiset of letters of $w$ is the \emph{content}. The history is \emph{complete} if $m_0$ is squarefree (equivalently, by \cref{thm:predecessors}(b), $m_0$ has no predecessor, so the history cannot be extended backwards).
\end{definition}
By \cref{lem:R} states increase strictly, so a history determines the set of its states, and $m_l=m_0\prod_j R(N_j)$ depends only on the source and the content.

\subsection{Vocabulary: state, history, content, endpoint, swap class, rewrite class}
The words are used in a fixed sense throughout.
\begin{itemize}[leftmargin=1.6em,itemsep=1pt]
\item A \emph{state} is an integer $m$. A \emph{history} is a path $(m_0;w)$. Its \emph{content} is the multiset of its labels and its \emph{endpoint} the last state. Two histories have the same endpoint from the same source iff they have the same multiplier $\prod R(N_j)$ (\cref{sec:endpoints}).
\item Two histories of the same source are \emph{swap-adjacent} if their words differ by exchanging two adjacent \emph{distinct} labels (\cref{def:swapgraph}); the connected components of the resulting graph are \emph{swap classes}. Swap classes lie inside content classes, which lie inside endpoint classes.
\item A \emph{block rewrite} replaces a contiguous block of a word by another word of equal multiplier, possibly with a different content (\cref{def:rewrite}); the components of swaps and block rewrites together are \emph{rewrite classes}.
\end{itemize}
\Cref{fig:objects} displays the nesting for the endpoint $72$. A function of the endpoint alone sees none of the finer relations; the manuscript's results are statements about the finer relations.

\subsection{Valuation walks}
For an integer $m$ write $u(m)=(v_p(m))_p$. A step with label $N$ requires $u_p\ge v_p(N)$ for all $p$ and changes $u$ by $\Delta(N)=(\Delta_p(N))_p$. A history is therefore a lattice walk in the nonnegative orthant with step set $\{\Delta(N)\}$: $p$ is a resource of capacity $u_p$, a label divisible by $p$ consumes it, and a label $N$ with $p\mid N+1$ produces it. The total displacement of a walk is determined by the content, so it is the \emph{route} of the walk---not its endpoint---that distinguishes histories. This is the picture behind \cref{thm:threshold} below; it is inherited arithmetic (record \texttt{D-W}).
\prov{CW-0337}{D-W}{PROVED}{FULL}

\begin{remark}[model boundaries]\label{rem:bounds}
Only integer states and the grammar of \cref{def:select} are used. Rational states, the collective map of the programme's FIELD-I, the occurrence geometry of rational states and the move complexes of rational states are not used; where an integer move complex enters (\cref{sec:geometry}) it is the complex of an integer endpoint, which is not a rational extension.
\end{remark}
